\begin{helper}[A component swap between two missing-color ends preserves non-elementarity]
Let \(X\) be a vertex set in a properly partially edge-colored graph.
Let \(u\ne v\) belong both to \(X\) and to one
\((\alpha,\beta)\)-component \(C\), where \(\alpha\ne\beta\), with
\(u\) missing \(\alpha\) and \(v\) missing \(\beta\).  If \(X\) is
not elementary before interchanging \(\alpha,\beta\) on \(C\), then it
is not elementary afterward.
\end{helper}
\begin{proof}
By \noderef{TwoColorComponentSpanningPath}, properness
\noderef{IsProperEdgeColoring} supplies a simple path whose support is
the component \noderef{TwoColorComponent} \(C\).  Since \(u\) misses
\(\alpha\) and \(v\) misses \(\beta\),
\noderef{MissingColorIncidenceEndpoint} makes each of their vertex-nodes
have degree at most one.  An internal support vertex of a simple path
has its two distinct path-neighbors, so neither \(u\) nor \(v\) is
internal.  They are distinct and belong to the support, and hence they
are exactly the two ends of this nontrivial spanning path.

We first classify the two affected missing colors at every vertex-node
of \(C\).  The unique path-neighbor of the end \(u\) is an edge-node
adjacent to \(u\),
so \noderef{TwoColorIncidenceGraph} says that it is an incident edge of
old color \(\alpha\) or \(\beta\).  Since \(u\) misses \(\alpha\), the
definition \noderef{MissingColors} excludes the first possibility.
Thus this edge has old color \(\beta\), and consequently \(u\) does not
also miss \(\beta\).  Symmetrically, the unique path-neighbor of \(v\)
is an edge-node of old color \(\alpha\), because \(v\) misses
\(\beta\); consequently \(v\) does not also miss \(\alpha\).

Every other vertex-node of \(C\) is internal on the spanning path and
has two distinct path-neighbors.  By
\noderef{TwoColorIncidenceGraph}, those neighbors are the identities of
two distinct incident edges, each colored \(\alpha\) or \(\beta\).
They cannot have the same color, because they share this vertex and
properness \noderef{IsProperEdgeColoring} forbids equal colors on
distinct adjacent edges.  Hence one has color \(\alpha\) and the other
has color \(\beta\), so an internal vertex misses neither affected
color.  We have proved that, among vertex-nodes of \(C\), only \(u\)
can miss \(\alpha\), and only \(v\) can miss \(\beta\), before the
swap.

Every edge-node on the spanning path belongs to \(C\), so
\noderef{TwoColorComponentSwap} reverses its affected color.  In
particular, the old \(\beta\)-edge at \(u\) becomes an \(\alpha\)-edge,
and the old \(\alpha\)-edge at \(v\) becomes a \(\beta\)-edge.  Together
with \noderef{TwoColorComponentSwapMissingBehavior}, this shows that
afterward \(u\) misses \(\beta\) but not \(\alpha\), while \(v\) misses
\(\alpha\) but not \(\beta\).  At each internal vertex, its two path
edges are merely exchanged, so it still misses neither affected color.
Vertices outside \(C\) keep their whole missing sets by
\noderef{TwoColorComponentSwapMissingBehavior}, and the same node says
that all colors outside \(\{\alpha,\beta\}\) keep their missing status
everywhere.

Because \(X\) is not elementary, unfolding
\noderef{ColoredVertexSetElementary} gives distinct \(x,y\in X\) and a
color \(c\) belonging to both sets from \noderef{MissingColors}.  If
\(c\notin\{\alpha,\beta\}\), the same witnesses work afterward.  Let
\(c=\alpha\).  Among vertices of \(C\), only the end \(u\) can miss
\(\alpha\) before the swap.  Thus either both witnesses lie outside
\(C\), in which case they remain witnesses, or exactly one is \(u\)
and the other, say \(x\), is outside \(C\).  In the latter case \(x\)
still misses \(\alpha\), while the other end \(v\in X\) now misses
\(\alpha\); also \(x\ne v\), since \(v\in C\).  Hence \(x,v\) witness
non-elementarity afterward.  The case \(c=\beta\) is symmetric, using
the new missing-\(\beta\) end \(u\).  In all cases two distinct members
of \(X\) share a missing color after the component swap, so \(X\) is
not elementary.
\end{proof}
